\chapter{量化 C++ 求职与入职}

\begin{introduction}
  \item 把语言知识组织成“定义、场景、代价、失败”四段式回答
  \item 用已运行的代码、测试、失败实验和 benchmark 回答代码题与系统设计题
  \item 把最终项目整理成可复现、不过度承诺的作品集证据
  \item 制定以代码库事实为起点的 30/60/90 天入职路线
\end{introduction}

\chaptermeta{已完成第 1--17 章；能从干净目录构建最终项目，解释四个公开 seam、失败边界与性能测量口径。}
{Python 经历不是需要隐藏的弱点：研究迭代、数据处理和验证习惯可以迁移；但 C++ 岗位会进一步追问对象生命周期、构建链接、内存模型、数值口径和成本证据。}
{最常见的失败不是“少背了一个术语”，而是把没有运行过的命令、没有测量的加速或没有参与的团队成果说成事实。本章只允许从可复现证据向外做有限推论。}

\section{先建立证据账本，再准备答案}

求职准备常从一份题库开始，结果是每道题都能说两句，却无法回答追问：“你在哪里用过？失败时看到了什么？为什么没有选另一种方案？”本书已经产生一条更可靠的证据链：源码说明做了什么，测试说明公开行为，失败实验说明边界，benchmark 说明怎样测量，README 说明怎样复现。面试答案只是这条链的索引。

先建立三列表，而不是先写简历形容词：

\begin{center}
\begin{tabularx}{\textwidth}{@{}p{0.20\textwidth}X p{0.27\textwidth}@{}}
\toprule
主张 & 可核验证据 & 必须同时声明的限制 \\
\midrule
能设计可测试组件 & 第 8、12、17 章的策略、成交/组合与端到端 seam & 教学项目不是生产撮合系统 \\
能诊断内存与并发错误 & ASan、UBSan、TSan 的隔离失败实验与稳定类别 & 工具只覆盖实际执行或探索到的路径 \\
能做性能工作 & checksum 门、原始样本、中位数/IQR、profiler 与 Amdahl 预算 & 本机观察不能外推到生产负载 \\
能连接 Python 与 C++ & 第 16 章的 dtype、连续性、owner 与 GIL 边界 & 核心可构建不等于已经交付 Python 包 \\
\bottomrule
\end{tabularx}
\end{center}

岗位名称并不统一，可先按工作证据分三类。研究工程更看重可复现数据流、数值口径和批量 Python 边界；通用 Quant Developer 还会深入领域建模、构建、测试、诊断和性能；低延迟或交易基础设施会继续追问布局、并发协议、延迟分布与网络/交易所语义。本书为前两类建立了主干，也提供第三类的基础证据，但没有网络栈、订单簿或生产延迟数据，不能把“学过缓存和线程”写成“有低延迟交易经验”。

\begin{pythonbridge}
Python 项目常用 notebook、图表和实验指标证明研究过程；C++ 项目还需要保留工具链、构建配置、目标图、测试退出码和 sanitizer 类别。两者都要求可复现，但证据颗粒度不同。
\end{pythonbridge}

立即练习：选择一种目标岗位，为自己写三条“主张—证据—限制”。证据必须能指向本书某条命令或源码；若只能写“熟悉”“了解”，先把它降级为学习目标。

\section{语言题：定义、场景、代价、失败}

语言题不是关键词接龙。稳定回答结构为：

\begin{enumerate}[label=\arabic*.]
  \item \textbf{定义：}它在 C++ 对象模型或编译模型中保证什么；
  \item \textbf{场景：}什么前置条件下使用，调用者能观察到什么；
  \item \textbf{代价：}复制、间接调用、共享状态、编译时间或接口约束是什么；
  \item \textbf{失败：}最小反例、诊断工具与修复方向是什么。
\end{enumerate}

以 RAII 为例，合格回答可以是：RAII 把资源的取得与一个对象的初始化绑定，把释放放进析构，因此所有正常离开作用域的路径都会按逆序清理；它适合文件、锁、动态存储和其他具有成对释放协议的资源；代价是必须明确所有权和对象寿命，移动或共享还会改变类型设计；若把 \texttt{span} 或裸指针保存到 owner 之外，RAII 只会准时释放资源，不能阻止悬空借用，第 6、13、16 章的 ASan 实验会报告 use-after-free。

这比“RAII 就是智能指针、可以避免泄漏”更准确。后一句漏掉栈对象、文件和锁，也掩盖了 \texttt{shared\_ptr} 环、错误 owner 和悬空观察者。Python 的 \texttt{with} 也表达作用域清理，但普通 Python 对象的析构时机不应被当作跨实现协议；C++ 自动对象的作用域退出和栈展开则是语言级清理边界。

同一结构可以压缩到其他高频问题：

\begin{center}
\begin{tabularx}{\textwidth}{@{}p{0.20\textwidth}X@{}}
\toprule
问题 & 回答骨架 \\
\midrule
按值还是 \texttt{const T\&} & 值拥有独立对象，适合需要保存/修改副本或小型可廉价复制值；引用只读借用且避免一次复制，但不延长寿命，保存引用会引入悬空风险。 \\
模板还是虚函数 & 模板在编译期对具体类型实例化，利于内联与约束诊断但增加编译/代码体积；虚接口支持运行期替换，承担间接调用、对象寿命与 ABI 约束。 \\
\texttt{vector} 何时失效 & 超过 capacity 的增长可搬迁存储，使旧指针、引用和迭代器失效；\texttt{reserve} 只降低搬迁概率，不允许无限增长，也不改变 size。 \\
异常还是返回值 & 正常缺席用 optional，需枚举的可处置失败用显式结果，跨层且无法局部处理的异常路径要满足保证并在掌握策略的边界记录一次。 \\
\bottomrule
\end{tabularx}
\end{center}

失败诊断：若回答只有定义，追问一个反例就会中断；若只有“更快”“更安全”，追问成本就会暴露没有边界。恢复方法不是继续堆术语，而是回到一个实际源文件、一个失败输入和一条工具输出。

立即练习：分别用四段式回答“\texttt{std::move} 做了什么”和“atomic 为什么不能自动给 Portfolio 一致快照”。每题限 90 秒，并必须引用第 6 或第 15 章的一个失败例子。

\section{把全书变成可检索的问题库}

零散复习很容易重复自己已经会的内容。更有效的方法是把问题按证据类型索引；回答卡只写关键词和源码位置，不抄整段标准定义。下面十二个入口覆盖了从语言到生产边界的主要追问：

\begin{center}
\begin{tabularx}{\textwidth}{@{}p{0.18\textwidth}p{0.35\textwidth}X@{}}
\toprule
主题 & 起始问题 & 实验证据 \\
\midrule
编译模型 & 声明、定义、编译与链接如何分工？ & ch01/ch03 的编译与链接失败 \\
类型与数值 & 窄化、整数除法和浮点容差为何不同？ & ch02 的编译诊断；ch16 的 oracle \\
借用与所有权 & 值、引用、指针和 owner 怎样选择？ & ch03/ch06；ASan use-after-free \\
容器与迭代器 & 为什么选 vector/map，何时失效？ & ch04/ch07 的失效和访问实验 \\
领域组件 & 不变量、optional 和虚接口怎样分工？ & ch08 的四个组件 seam \\
泛型 & 模板、concept 与运行期多态怎样权衡？ & ch09 的两类编译诊断 \\
错误语义 & absence、坏输入、异常和日志何时使用？ & ch10 的恢复管道与一次日志 \\
构建 & target、usage requirement 和配置是什么？ & ch11 clean build 与链接失败 \\
质量工具 & 测试、警告、静态分析和 sanitizer 各看到什么？ & ch12 的隔离故障矩阵 \\
性能 & 布局、分配、benchmark 与 profiler 如何形成证据链？ & ch13--14 checksum 与样本 \\
并发与边界 & happens-before、背压、GIL 与 owner 如何闭合？ & ch15--16 的失败实验 \\
项目 & 成本、估值、订单生命周期和生产边界是什么？ & ch17 的 CLI、失败与 benchmark \\
\bottomrule
\end{tabularx}
\end{center}

练习不是把每题写成文章，而是做三轮不同约束。第一轮每题 30 秒，只给定义和一个证据位置；第二轮 90 秒补场景、代价与失败；第三轮让同伴从任意词追问“为什么、反例、怎么测”。若回答中出现没见过的语法或平台结论，标成缺口，不在现场编造。

例如“\texttt{shared\_ptr} 是否线程安全”至少要拆成两层：不同 shared pointer 对象对同一控制块的所有权计数操作可以安全并发，但被指对象的业务状态不会因此自动同步；多个线程修改 Portfolio 仍需第 15 章的复合不变量协议。只回答“是”或“否”都丢失了对象边界。

立即练习：从表中随机抽三题，每题画出“公开保证—反例—工具—项目位置”四格卡。第二天不看原文重答；若只能背出工具名，回到失败源实际运行一次。

\section{代码题：先冻结契约，再写循环}

量化代码题常披着收益率、窗口或订单的外衣，但评分点仍是边界、数据结构、复杂度和可验证性。拿到题目后按以下顺序工作：复述输入输出；确认空输入、非法值、单位与时间顺序；写两个手算例；选择状态与不变量；先写清楚版本；最后分析复杂度、溢出和测试。不要一开始就猜面试官想要的“最优技巧”。

以“给定按时间排列的权益，返回最大回撤”为例，先问：峰值是否包含初始现金？权益能否非正？返回比例还是金额？空序列返回零还是错误？本项目的契约是从 initial cash 开始维护历史峰值，只在峰值为正时计算比例，并把空曲线的摘要留为零值。

独立手算序列 \(100,80,120,90\)：先有 20\% 回撤，随后峰值更新为 120，最后回撤 25\%，答案是 25\%。一次扫描只需保存 \texttt{peak} 和 \texttt{max\_drawdown}，时间 \(O(n)\)、额外空间 \(O(1)\)。实际实现位于 \path{project/include/quant/statistics.hpp}，端到端行为由 \path{project/tests/test_backtest.cpp} 约束。

常见错误是先求整个序列的最大值 120，再把它用于更早的 80；这会制造 33.3\% 的“未来峰值回撤”，属于 look-ahead。另一个错误是测试用同一循环重新计算 expected；正确 oracle 是上面的手算 25\%。Python 的 \texttt{Series.cummax} 能表达相同历史状态，但面试中仍要说明其时间语义、缺失值政策和物化成本，不能用 API 名称替代契约。

若时间允许，再补性质：单调递增序列回撤为零；添加一个新高点本身不会增加回撤；把所有正权益乘以同一正数，比例回撤不变。这些性质补充例子，却不能替代具体边界测试。

立即练习：对 \(100,110,88\) 先手算 20\%，再口述一次扫描中的每次状态变化；然后解释为什么用最终最大值计算所有历史点会发生未来信息泄漏。

\section{限时作业与现场代码评审}

take-home 或现场扩展题通常同时考察实现与交付。拿到题目后先写一段 decision log：已确认的输入输出、允许的依赖、时间预算、明确非目标和验收命令。若题目有歧义，保存问题与采用的假设；沉默地猜中一个版本不如让审阅者看见决策边界。

实现仍使用纵向 tracer bullet。先让一个代表性输入从 CLI 或公开函数得到手算结果，再补一个稳定失败；每轮保持可构建。提交前从新目录运行格式化、警告、测试和示例，不把本机 build 目录、绝对路径或依赖缓存当作项目的一部分。README 只写实际执行过的命令，并说明编译器与平台。

评审时不要逐行朗读代码。先用两分钟说明契约、所有权和数据流，再主动指出最可能出错的边界、测试如何对其敏感以及时间不足时没有做什么。常见追问包括：为什么不是另一种容器；哪个引用可能悬空；异常后公开状态怎样；整数在哪里可能溢出；线程安全是对象级还是操作级；复杂度中的 \(n\) 是什么；benchmark 是否包含 I/O。

可以用最终项目做一次模拟评审。固定点选在加入成本配置之前，diff 只展示 ExecutionConfig 如何沿 BacktestConfig 进入撮合、Fill 和 Portfolio；测试的独立账本先红再绿。审阅者若问“为什么不同时支持阶梯手续费和真实盘口”，回答当前 tracer bullet 只证明固定费用与确定性滑点，新增模型必须有新 oracle，不能把未来扩展塞进当前接口。

失败诊断：一次性提交几百行、测试只断言“没有异常”、expected 由程序自动重写、README 命令依赖旧 build、为了显得完整引入未使用框架，都会降低可审阅性。修复是缩小 diff、恢复独立值、删除 speculative generality，并让陌生人从 clean build 复现。

立即练习：假设只有 90 分钟为项目增加“非法负费用返回状态 2”。写出 10 分钟澄清、15 分钟红灯、30 分钟实现、15 分钟回归、10 分钟 README、10 分钟自评的时间盒；实际执行后比较偏差，而不是事后声称按计划完成。

\section{性能题：数字之前必须有正确性门}

被问“为什么用 C++ 会更快”时，不要把语言名称当作 benchmark。Python 循环可能承担动态分派和对象开销，NumPy 运算也可能已经在优化后的 C/C++ 内核中执行；C++ 能提供连续布局、可控分配、静态优化和较低的语言边界开销，但实际收益取决于热点、数据移动和端到端占比。

一个可审阅的性能回答依次包含六项：业务 oracle 一致；代表性工作量；CPU、编译器和 Release 配置；预热与原始样本；中心、离散或尾部统计；测到与未测到的边界。第 17 章 benchmark 的输出结构正好可作为答题证据：

\begin{lstlisting}[numbers=none]
benchmark-ok events=200000 samples=11 warmups=2 equity=10100 fills=1 raw-us=<11 samples> median-us=<value> iqr-us=<value> checksum=131300
\end{lstlisting}

这里 \texttt{equity}、\texttt{fills} 与 \texttt{checksum} 是正确性门；\texttt{raw-us} 是可复核观察；median/IQR 是摘要。CTest 调用 \path{tools/validate_ch17_benchmark.py}，确认恰有 11 个非负整数样本并从原样本复算统计。时间本身随机器变化，所以不写入跨机器 expected。

若一次优化让 40\% 的端到端时间缩短四倍，总加速只有
\[
\frac{1}{0.6+0.4/4}\approx1.43,
\]
而不是四倍。回答还应说明 profiler 如何得到 40\%，以及改动是否增加内存、复杂度或尾部延迟。只展示最快一次、使用 Debug、比较不同 checksum、删除输入验证或不记录环境，都会让数字失去归因能力。

\begin{pythonbridge}
Python profiler 可定位解释器或调用边界热点，C++ profiler 可进一步看到函数、分配与硬件事件；跨语言系统必须同时观察批大小、复制、dtype 转换和 GIL 区间。只测一个空 pybind11 调用，不能代表真实批量接口。
\end{pythonbridge}

立即练习：把“C++ 版本快 3 倍”改写成一条可证伪结论，至少补上 workload、环境、oracle、原始样本统计和未覆盖边界；缺一项就标记为待测，而不是自行补数字。

\section{系统设计题：从数据契约画到失败边界}

考虑题目：“设计一个读取历史行情、运行策略并输出可审计绩效的回测服务。”先确认范围：输入是批文件还是实时流？需要哪些资产、币种和公司行动？同一输入是否必须确定性重放？吞吐、完成时限、并行任务数和审计保留期是什么？若需求未给出，不要直接画 Kafka、数据库和微服务。

最小单机数据流可以复用第 17 章：
\[
\text{validated data}\rightarrow\text{strategy intent}\rightarrow
\text{execution fact}\rightarrow\text{portfolio ledger}\rightarrow
\text{versioned result}.
\]
每条箭头都要回答所有权和失败：输入解析拒绝哪类坏行；策略能否修改账户；成交如何保存费用、滑点和订单 ID；组合怎样保证原子账本；统计需要哪一时刻的完整价格；结果怎样关联代码版本、数据版本、配置和运行 ID。

扩到多资产时，不能只把 \texttt{vector} 里的代码放宽。估值需要 \texttt{Symbol -> Price} 的时间一致映射、缺失/陈旧价格政策和币种转换；事件乱序需要 watermark 或明确排序；部分成交、撤单与拒绝需要状态机和幂等事件 ID。当前项目选择稳定拒绝混合代码：

\begin{lstlisting}[numbers=none]
model-boundary-error single-symbol backtest requires one symbol; define a complete price map before enabling multi-asset valuation
\end{lstlisting}

这条状态 2 失败由 \path{project/failures/mixed_symbols.cpp} 与 CTest 固定。面试时承认边界，比用最后一个价格算出漂亮但错误的总权益更可信。若要演化，第一条 tracer bullet 应是两资产、同一估值时刻、完整价格映射和独立手算账本；随后才加入 stale price、公司行动和并行调度。

并发扩展也从协议开始。可以按相互独立的策略/参数任务分区，每个任务拥有自己的账本；共享队列必须说明容量、满策略、close、stop、排空和异常传播。单次回测内部若共享可变 Portfolio，就要证明复合现金/仓位不变量的 happens-before，而不是把两个字段分别改成 atomic。

可观测性至少记录输入数据版本、策略/二进制版本、配置哈希、开始结束时间、拒绝计数、处理事件数和结果校验；性能指标必须与正确性和业务结果关联。Python 可以承担研究编排、参数生成和结果分析，C++ 核心通过批量、明确 owner 的边界承接热点；逐事件跨语言调用通常会让边界成本吞掉内核收益。

立即练习：在纸上给上述系统增加一个“MSFT 在估值时缺价”的输入。画出数据从读取到结果的路径，明确选择拒绝、最近价加 stale age 或不完整结果三者之一，并写出调用者可观察的状态，不能默认为零。

\section{完整模拟：从代码题追问到系统设计}

下面是一条可以独立排练的追问链。不要先读示范要点；开 35 分钟计时器，拿纸、编辑器和一个空终端完成，再按每步证据自评。

\subsection{第一轮：最大回撤的流式接口}

面试官先给固定序列题，随后追问：“权益逐项到达，不能保存全部历史，接口怎样改？”候选人应指出算法本来只需要历史 peak 与 maximum drawdown，因此可把状态放进一个小对象，每次 \texttt{observe(equity)} 更新，\texttt{value()} 返回当前结果；状态属于这个计算对象，调用者拥有其寿命。还要确认初始 cash 如何进入、非有限值怎样处理、是否允许跨运行复用。

若进一步要求多线程同时调用，不应立刻加 atomic。先问是否能按回测 run 分区，让每个 run 独占状态；这是最简单且最容易确定性重放的方案。只有确实共享一个序列时才讨论锁和事件顺序，而“同时到达”的权益如果没有全序，最大回撤本身就缺少定义。

白板验收至少包含：空输入；单个值；单调上升；\(100,80,120,90\) 得 25\%；\(100,110,88\) 得 20\%；非法 NaN 的显式政策。测试 expected 来自手算，不能调用另一个相同循环生成。

\subsection{第二轮：把算法放回回测系统}

接着追问：“这个对象接在系统哪里？”它应只读取按契约生成的权益快照，不修改 Portfolio，也不直接解析 CSV。Portfolio 应用 Fill 后产生某估值时刻的 snapshot，统计阶段观察 snapshot；这样策略、撮合和统计可以分别测试。若统计对象直接读取 Portfolio 私有 map，它就跨过公开 seam，并使容器重构破坏测试。

再追问：“两种结果不同怎样排查？”先固定相同事件、费用、滑点、估值价格和浮点口径，比较第一处权益曲线分歧；用手算小序列定位是账本还是统计。不能先调容差或只比较最终收益，因为中间一次错误可能被后续价格抵消。

\subsection{第三轮：服务化与规模}

最后给一个明确的练习假设：历史任务可以彼此独立，每个任务输入一个已版本化数据集与策略配置，要求失败可重放，不要求实时撮合。此时可按 run ID 分区工作队列；worker 内保持单所有者账本；结果存输入哈希、二进制/策略版本、配置、诊断和摘要。队列满时选择阻塞提交或显式拒绝并重试，不能无限增长。

若面试官把需求改成“共享实时行情供多个策略消费”，新的核心问题是事件顺序、慢消费者、数据可丢性和恢复，而不是把原架构名称换成流平台。候选人应询问是否允许丢行情、如何快照恢复、每个策略是否需要完全相同序列，以及关闭时要排空到哪个 watermark。

\subsection{自评量表}

每项 0--2 分：是否先澄清契约；是否给独立手算；是否说明所有权；是否区分正常缺席与失败；是否给时间/空间复杂度；是否在性能前守住 oracle；是否声明当前单代码限制；是否给测试、诊断和恢复。12 分以下回到相应章节补一次运行证据，不能只重背示范回答。

常见失败有三类：写出循环却不定义时间语义；画出服务却不说明账本唯一写入者；给出并发优化却没有 workload 和停止协议。修复顺序也是契约、所有权、正确性、可观测性、测量，最后才是组件选型。

立即练习：把练习假设改成“允许取消正在运行的任务”。补出 created、running、stop-requested、completed、failed 与 cancelled 的合法转移，说明重复取消、结果提交和 worker 异常的幂等规则。

\section{作品集：让陌生人从空目录复现}

作品集项目的第一位读者不是面试官，而是未来的自己或一个不了解上下文的审阅者。\path{project/README.md} 是本书最终项目的证据入口，包含范围、Linux/WSL clean build、两条精确摘要、11 项 CTest、benchmark 口径、一次失败修复和明确非目标。最小复现为：

\begin{lstlisting}[language=bash,numbers=none]
cmake -S project -B build/ch17 -DCMAKE_BUILD_TYPE=Release
cmake --build build/ch17 -j
ctest --test-dir build/ch17 --output-on-failure
./build/ch17/quant_backtest project/data/sample.csv 1.0 10.0
\end{lstlisting}

合格的五分钟项目叙事按“问题—边界—证据—失败—限制—下一步”展开：Python 原型容易把策略和账户混在一起；项目把事件、意图、成交和账本分开；四个 seam 与手算成本账本守住行为；混合代码暴露估值契约缺失，于是先增加状态 2 拒绝；当前仍是同步、单代码教学引擎；下一步先做完整价格映射，不先做无锁队列。

简历 bullet 也只压缩这个事实。例如：“用 C++20 实现可独立构建的事件驱动回测教学项目，以四个公开行为 seam 验证 CSV、策略、成交/组合和端到端账本；加入费用/滑点配置、混合代码显式拒绝及正确性门控的 11 样本 benchmark。”只有在自己确实完成构建、测试和解释后才能使用。不要改写成“构建生产级高频平台”或“收益提高 20\%”。

一次可信的失败修复需要保存：失败输入、期望与实际、退出码/诊断类别、根因、最小改动和回归。混合代码故事的根因不是“vector 有 bug”，而是 Portfolio 虽存多代码仓位，最终 snapshot 却只有一个 mark price；修复也不是伪造总权益，而是缩小公开契约并登记未来所需价格映射。

立即练习：让同伴只看 README，从新目录复现两条摘要。你不能口头补步骤；若失败，把缺失命令、依赖或限制写回 README。然后用 90 秒讲清项目，录音中每个数字都必须能在输出或测试中找到。

\section{把能力缺口变成下一项证据}

诚实声明限制不是结束语，还要把它转换成最小学习项目。缺口描述必须是可验证行为，而不是“继续深入 C++”。

\begin{center}
\begin{tabularx}{\textwidth}{@{}p{0.22\textwidth}X p{0.28\textwidth}@{}}
\toprule
目标缺口 & 下一项 tracer bullet & 不应提前声称 \\
\midrule
多资产估值 & 两资产、同一时刻完整价格映射；缺价状态 2；手算总权益 & 已支持组合级历史回测 \\
订单生命周期 & 稳定 order/event ID；部分成交后撤单；重复 Fill 幂等 & 已模拟真实交易所 \\
Python 交付 & 一个批量 NumPy 边界；dtype/stride/owner 失败；打包与导入测试 & 已有成熟 Python SDK \\
低延迟基础 & 固定 workload 的分阶段 latency 分布；profiler/硬件计数；业务 checksum & 达到生产微秒延迟 \\
真实数据语义 & 数据版本、交易日历、复权/公司行动的一条端到端样例 & 策略结果可直接用于实盘 \\
团队协作 & 一个真实 issue、评审反馈、修订与合并记录 & 有大型团队领导经验 \\
\bottomrule
\end{tabularx}
\end{center}

每个 tracer bullet 先写完成判据。例如多资产项不是“把单代码检查删掉”，而是：AAPL 与 MSFT 各有非零仓位；估值时两者价格齐全则精确匹配手算；缺任一价格稳定失败；旧单代码与成本 seam 不变。这样新的简历材料只能在事实完成后自然升级。

选择下一项时看目标岗位和证据密度。研究工程岗位可能优先做数据版本与批量 Python 边界；交易基础设施岗位可能优先做订单状态机、重放和并发协议；通用 Quant Developer 可以先补多资产与错误可观测性。不要同时开启三条横向重写。

立即练习：从表中选择一项，写出一个红灯输入、一个独立 oracle、最小公开接口、明确非目标和完成后才能使用的一条项目描述。若无法在一周左右得到可观察结果，把切片继续缩小。

\section{行为面试：只讲亲手完成的事实}

STAR 结构有用，但容易把“团队背景”写成自己的成果。更安全的证据卡包含：情境与约束；你实际负责的动作；动作前后的可观察结果；保留的代码/测试/评审；仍未解决的限制。若项目是个人学习项目，就明确说个人项目，不虚构客户、同事或线上事故。

本书提供三个可转化素材，但只有亲手复现相应步骤后才能使用：成本配置的红—绿循环可回答如何小步交付；checksum mismatch 可回答如何拒绝无效性能结论；混合代码状态 2 可回答如何在模型不完整时避免静默错误。结果不是“帮助公司节省资金”，而是具体测试从红到绿、错误进程变成稳定非零退出、旧 seam 继续通过。

面对“说一次你犯的错误”，可以说明一开始把性能或功能目标写得过宽，随后用独立 oracle 或失败输入发现问题，缩小改动并补回归。不要把“我太追求完美”当作错误，也不要声称自己处理了没有证据的生产事故。面对团队冲突题，个人项目不能生成团队故事；应换用真实课程、实习、开源评审或工作经历，并保留同样的事实边界。

立即练习：为一件真实经历写六行证据卡，圈出每个“我”动词，并删除无法由提交、测试、文档、评审或真实见证支持的结果。若只剩学习项目，也可以诚实回答“这是个人项目，协作方面我用另一段经历说明”。

\section{入职前三个月：先降低未知，再交付改动}

入职计划不是“第 90 天重构架构”。量化 C++ 系统往往同时包含数据口径、交易语义、构建环境、历史兼容和性能约束；最初目标是把未知变成可验证地图。

\begin{center}
\begin{tabularx}{\textwidth}{@{}p{0.13\textwidth}p{0.25\textwidth}X@{}}
\toprule
阶段 & 核心问题 & 可交付证据 \\
\midrule
0--30 天 & 怎样可靠构建、测试、运行和观察一条真实路径？ & 开发环境复现记录；目标/测试图；一条数据到结果的调用与所有权图；值班、发布和回滚入口；术语与负责人清单。 \\
31--60 天 & 关键正确性、数据和性能口径是什么？ & 一个业务账本手算；错误分类与日志/指标路径；带环境的基线；一次小缺陷复现或文档修复；评审反馈记录。 \\
61--90 天 & 哪个最小改动能安全进入真实流程？ & 明确验收的 tracer bullet；红—绿证据；兼容/数据/性能影响；发布与回滚方案；上线后观察及后续边界。 \\
\bottomrule
\end{tabularx}
\end{center}

前 30 天先使用团队规定的编译器、依赖、数据和测试入口，不自行建立平行脚本；读一条真实业务路径时同时记录对象所有者、线程、错误和数据版本。遇到慢测试或 flaky benchmark，先复现其环境与统计口径，不凭一次本机结果删除它。

31--60 天选择一个范围小但能穿过真实 seam 的问题，例如补充一个稳定错误上下文、验证一个数据边界或修正文档化构建缺口。性能工作先取得 profiler 和业务 oracle；数值差异先确认单位、舍入和 reference。主动请领域专家审阅金融语义，请平台专家审阅发布与可观测性。

61--90 天的首次代码改动应可回滚、可观察、兼容已有调用者，且失败半径有限。提交说明包含为什么改、什么不改、测试证据、性能/数据影响和回滚。不要把“入职展示能力”变成一次大重写；能在旧系统约束下安全交付，比展示新语法更有价值。

第一周需要主动提问，但问题应指向可操作事实。构建方面问“CI 的权威编译器、依赖锁定和 clean build 命令是什么”；数据方面问“时间、币种、复权和缺失值的标准口径在哪里”；运行方面问“如何关联输入版本、二进制和结果”；故障方面问“哪些错误可重试、谁决定降级、怎样回滚”；性能方面问“哪条业务 oracle 必须不变、基线环境与尾部指标是什么”；协作方面问“哪些目录由谁负责、评审和发布需要哪些批准”。

这些问题的答案应进入个人代码库地图，而不是只留在聊天记录。对尚未确认的事项标注 owner 和下一步；不要把同事的临时猜测写成团队契约。读到历史 workaround 时，先找关联 issue、测试和事故背景，再判断它是否仍有必要。

入职早期的红旗包括：只有某位同事机器能构建；Release 与测试配置不同却无人说明；测试依赖共享未版本化数据；错误被日志吞掉后仍返回成功；性能报告没有业务 checksum；关键对象跨线程但 owner 不明。发现红旗先建立最小复现和影响范围，再通过团队流程报告，不能直接删除“奇怪代码”。

\begin{pythonbridge}
Python 团队也需要同样的代码库地图和安全改动；C++ 环境还要额外确认 ABI、编译选项、链接依赖、对象寿命、线程模型和 Release/Debug 差异。不要因自己熟悉 Python 就绕过生产 C++ 的标准入口。
\end{pythonbridge}

立即练习：把“第一个月熟悉代码、第二个月做需求、第三个月优化性能”改写成上表形式，每阶段至少给一个文件、命令、图、测试或可观测结果，并列出需要向同事确认而不能自行假设的事项。

\begin{interviewcheck}
进行一次 45 分钟模拟：5 分钟项目叙事；8 分钟回答 RAII、vector 失效、模板/虚函数与 atomic 快照；12 分钟完成最大回撤契约和算法；15 分钟扩展多资产回测设计；5 分钟说明入职首个安全改动。记录每次没有证据、没有边界或无法复算的回答，回到对应章节补实验，而不是只润色措辞。
\end{interviewcheck}

\section{输出任务与自测}

输出任务：建立一个求职证据包，至少包含五项。

\begin{enumerate}
  \item 一页目标岗位表：三条主张—证据—限制和尚缺能力；
  \item 八道语言题的四段式回答，每题引用一个源码或失败实验；
  \item 最大回撤代码题的契约、两个手算例、复杂度与未来信息失败例；
  \item 一页回测系统设计和五分钟项目讲稿，附 clean build、11 项 CTest、两条精确摘要及一次失败修复；
  \item 一页 30/60/90 天计划，每阶段都有可观察交付物、需确认事项和风险边界。
\end{enumerate}

验收时由另一人只使用证据包和仓库复现。所有命令必须从空构建目录成功；所有数字能指向输出；性能结论带环境、原始样本摘要和限制；行为素材不把书中作者或虚构团队的工作算给自己。若没有同伴，就隔一天在新终端按文档逐条执行并保存退出码。

自测：
\begin{enumerate}
  \item 四段式语言回答分别解决什么追问？只给定义缺少什么？
  \item RAII 为什么不能自动阻止借用对象悬空？Python \texttt{with} 与它哪里相似、哪里不同？
  \item 最大回撤为什么必须使用历史峰值，不能使用最终全局最大值？
  \item “C++ 快三倍”至少缺少哪五类证据？
  \item 多资产回测从单代码演化时，价格映射还必须包含哪些时间与缺失语义？
  \item 为什么把两个 Portfolio 字段分别改成 atomic 仍不能得到一致快照？
  \item 五分钟项目叙事中的失败修复为什么必须包含原输入、根因和回归？
  \item 哪些表述会把教学项目夸大成生产交易平台？怎样改写？
  \item 个人项目能否回答团队冲突题？证据不足时应怎样处理？
  \item 首个入职改动为什么应优先小、可回滚和可观察，而不是架构重写？
\end{enumerate}

\section{本章小结}

求职与入职不是技术章节之外的包装。可信表达仍遵守本书的工程方法：冻结契约，用独立事实验证结果，保存失败边界，在测量前守住正确性，并明确没有完成什么。语言回答索引对象模型，代码题索引不变量，系统设计索引所有权与失败，作品集索引可复现命令，入职计划索引安全交付。能从证据出发并在追问中守住边界，才是这套 C++ 能力真正可迁移的完成标志。
